%
% crosshost_hubid_preemption_name_keyed_buggy.tex
%
% FIDELITY CONTROL for crosshost_hubid_preemption.tex (W12, lux-833p).
% This is the PRE-W11, BUGGY design: single-owner preemption keys on the
% declared ConnectMessage.name, not on HubId.  It exists to prove the
% correct model detects the exact defect W11 (lux-u9gb) fixed:
%   - Invariant 1 (at most one live per HubId) is VIOLATED here (goal FOUND):
%     two live connections sharing a HubId but declaring different names
%     both survive -- the double-owner state.
%   - Witness W2 (same-name, different-HubId coexistence) is DESTROYED here
%     (goal NOT found): a same-name arrival wrongly evicts a live, unrelated
%     Hub.
% Do NOT treat this as a description of shipped behaviour; the authoritative
% model is crosshost_hubid_preemption.tex.  See its Fidelity section.
%
% Models the code in:
%   src/punt_lux/display/client_registry.py   -- ClientRegistry.hub_fd_for
%   src/punt_lux/display/hub_reconciliation.py -- HubReconciliation._preempt_stale_hub
% and the invariant stated in
%   docs/architecture/system.tex "Cross-Host Transport and Trust",
%   subsection "Invariants", item 2.
%
% This document supersedes the name-keyed preemption clauses of
% hub_display_reconciliation.tex (its own SUPERSEDED banner points here):
% HI4 there ("preemption scopes to the shared name") described the pre-W11
% code and is corrected below.
%
% Type-check:  fuzz docs/crosshost_hubid_preemption_name_keyed_buggy.tex
% Model-check (goal predicates in the Verification section):
%   probcli docs/crosshost_hubid_preemption_name_keyed_buggy.tex \
%       -model_check -p DEFAULT_SETSIZE 3
%
% Formatting follows the fuzz house rules: \quad~ for continuation indent,
% schema lines under 80 columns, predicates broken at \land/\lor/\implies.
%

\documentclass[a4paper,10pt,fleqn]{article}
\usepackage[margin=1in]{geometry}
\usepackage{fuzz}
\usepackage[colorlinks=true,linkcolor=blue,citecolor=blue,urlcolor=blue]{hyperref}

\begin{document}

\title{Cross-Host Single-Owner Preemption:
a Fidelity Control (name-keyed, BUGGY)}
\author{W12 pre-W11 variant --- preemption keyed on \texttt{name}, which
violates DES-090 Invariant 2}
\date{September 2026}
\maketitle

\noindent\fbox{\parbox{\dimexpr\textwidth-2\fboxsep-2\fboxrule}{%
\textbf{Fidelity control --- BUGGY by construction.} This is the pre-W11
design, kept only to prove that
\texttt{crosshost\_hubid\_preemption.tex} detects the defect W11 fixed.
Preemption here keys on the declared \texttt{name}, not on \texttt{HubId}.
Invariant 1's negation is \emph{reachable} here (two live connections own
one \texttt{HubId}), and witness W2 (same-name, different-\texttt{HubId}
coexistence) is \emph{unreachable}. The authoritative model is
\texttt{crosshost\_hubid\_preemption.tex}.}}
\vspace{1em}

% ============================================================
\section{Introduction}
% ============================================================

The Display may aggregate Hub connections from several machines at once.
Each such connection carries, in its \texttt{ConnectMessage}, two
independent identifiers: a \texttt{HubId} --- a per-process identity, the
peer's hostname (cross-host, certificate-verified) paired with its
process id --- and a \texttt{name}, a human label that many Hub processes
may legitimately share (every production Hub today declares the identical
hardcoded name).  The two are not the same kind of thing, and the whole
point of this model is that single-owner preemption must key on the
former, never the latter.

DES-068 established, for the then-single same-host Hub, that a fresh
\texttt{kind="hub"} identify forcibly disconnects any predecessor
declaring the \emph{same identity} before recording itself --- so a dead
process's stale connection cannot linger and swallow clicks.  With one
Hub, ``same identity'' was unambiguous.  Cross-host makes two
independently-legitimate, simultaneously-live Hubs the ordinary case, and
the question ``which predecessor does a new connection preempt?'' now has
a right answer and a wrong one:

\begin{description}
  \item[right --- key on \texttt{HubId}.] A reconnect under the same
    \texttt{HubId} preempts its own stale predecessor; two connections
    with \emph{different} \texttt{HubId}s never preempt each other, even
    when they share a \texttt{name}.  This is the W11 (\texttt{lux-u9gb})
    change: \texttt{ClientRegistry.hub\_fd\_for} keys the lookup on
    \texttt{HubId}.
  \item[wrong --- key on \texttt{name}.] The pre-W11 code keyed the
    lookup on \texttt{ConnectMessage.name}.  Two Hubs sharing a name on
    different machines would then evict one another (a live, unrelated
    Hub wrongly disconnected), while two connections sharing a
    \texttt{HubId} but declaring different names would \emph{fail} to
    collapse --- both left live for one identity.
\end{description}

Two properties must hold across every interleaving of $N$ Hubs
connecting, reconnecting, and disconnecting:

\begin{description}
  \item[I1 --- at most one live owner per \texttt{HubId}.] No reachable
    state holds two distinct live connections declaring the same
    \texttt{HubId}.  This is Invariant 2 of the design, stated exactly.
  \item[I2 --- deadlock freedom.] No interleaving traps the model; some
    operation is always enabled.
\end{description}

Alongside these, two \emph{witnesses} must be reachable, or the
invariants would hold only because the mechanism they guard is dead
(Section~\ref{sec:verify}): different-\texttt{HubId} connections
\emph{coexisting} (preemption is not over-eager), and a same-name pair of
them coexisting in particular (the coexistence the wrong, name-keyed
design destroys).

% ============================================================
\section{Basic Types}
% ============================================================

\begin{zed}
  [CONN, HID, NAME]
\end{zed}

\texttt{CONN} is the set of Hub-Display connections (file descriptors)
the Display can hold open; three suffice to model two connections
contending over one identity plus a third unrelated one.  \texttt{HID} is
the set of \texttt{HubId} values a connection may declare --- two suffice
to exhibit both a same-identity reconnect and two distinct coexisting
identities.  \texttt{NAME} is the set of declared human labels; two
suffice to exhibit the (same \texttt{HubId}, different \texttt{name}) pair
the name-keyed variant mis-handles and the (different \texttt{HubId},
same \texttt{name}) pair it wrongly evicts.  All three carriers are
bounded by ProB's \texttt{DEFAULT\_SETSIZE}.

% ============================================================
\section{State}
% ============================================================

The state is flat: one schema.  A connection is ``live'' once it has
identified as \texttt{kind="hub"} and its socket is still open --- the
state in which it owns content and services clicks.  \texttt{hubOf} and
\texttt{nameOf} record what each live connection declared.

\begin{schema}{Owners}
  live : \power CONN \\
  hubOf : CONN \pfun HID \\
  nameOf : CONN \pfun NAME
\where
  \dom hubOf = live \\
  \dom nameOf = live
\end{schema}

The two domain equalities are structural, not the safety property: a live
connection has, by construction, declared exactly one \texttt{HubId} and
one \texttt{name} (both arrive in its single \texttt{ConnectMessage}), and
a departed one has neither.  ``At most one live per \texttt{HubId}'' ---
the injectivity of \texttt{hubOf} on \texttt{live} --- is I1, deliberately
\emph{not} written here: placing it in the state schema would enforce it
as a guard on every successor and silently disable any operation that
could break it, masking the very defect the name-keyed variant
(Section~\ref{sec:fidelity}) must be able to reach.

% ============================================================
\section{Initialization}
% ============================================================

\begin{schema}{Init}
  Owners'
\where
  live' = \emptyset \\
  hubOf' = \emptyset \\
  nameOf' = \emptyset
\end{schema}

A single initial state: no connection is live, and nothing has been
declared --- the cold start before any Hub has connected.

% ============================================================
\section{Operations}
% ============================================================

\subsection{Identify --- single-owner preemption, keyed on \texttt{HubId}}

A fresh connection \texttt{c?} identifies as \texttt{kind="hub"} declaring
\texttt{HubId} \texttt{h?} and name \texttt{n?}.  \texttt{c?} must not
already be live: a reconnect always arrives on a new file descriptor, so
the predecessor is a \emph{different} \texttt{CONN} value, and idempotent
re-identification of an already-live fd is not a distinct transition here.

The preempted set is exactly the live connections declaring the same
\texttt{HubId} \texttt{h?} --- \texttt{hub\_fd\_for(h?)} in the code,
lifted from ``the one such fd'' to ``every such fd'' so the model can
never assume uniqueness it is trying to prove.  Each is dropped
(\texttt{remove\_client}); \texttt{c?} is then recorded as the sole live
owner of \texttt{h?}.

\begin{schema}{Identify}
  \Delta Owners \\
  c? : CONN \\
  h? : HID \\
  n? : NAME
\where
  c? \notin live \\
  live' = (live \setminus \{ c : live | nameOf~c = n? \}) \cup \{ c? \} \\
  hubOf' = (\{ c : live | nameOf~c = n? \} \ndres hubOf) \oplus \{ c? \mapsto h? \} \\
  nameOf' = (\{ c : live | nameOf~c = n? \} \ndres nameOf) \oplus \{ c? \mapsto n? \}
\end{schema}

\subsection{Disconnect --- an ordinary drop}

A live connection's socket closes on its own --- a network blip, a
graceful shutdown, or a process exit.  \texttt{c?} leaves \texttt{live}
and its declarations are forgotten.  No new mechanism: a dropped
cross-host connection is, to the socket listener, an ordinary fd closing,
handled by the existing per-fd reaping the same-host case already uses.

\begin{schema}{Disconnect}
  \Delta Owners \\
  c? : CONN
\where
  c? \in live \\
  live' = live \setminus \{ c? \} \\
  hubOf' = \{ c? \} \ndres hubOf \\
  nameOf' = \{ c? \} \ndres nameOf
\end{schema}

% ============================================================
\section{Invariants}
\label{sec:inv}
% ============================================================

\paragraph{Invariant 1 --- at most one live owner per \texttt{HubId}.}
\texttt{hubOf} is injective on the live set:
\[
  \forall c_1, c_2 : live @ \\
  \quad~ hubOf~c_1 = hubOf~c_2 \implies c_1 = c_2.
\]
\texttt{Identify} removes \emph{every} live connection declaring
\texttt{h?} before adding \texttt{c? \mapsto h?}, so no two live
connections can share a \texttt{HubId}; \texttt{Disconnect} only shrinks
\texttt{live}, which cannot create a collision.  The check
(Section~\ref{sec:verify}) asks whether the negation --- two distinct live
connections with equal \texttt{hubOf} --- is reachable.

\paragraph{Invariant 2 --- deadlock freedom.}
\texttt{Disconnect} is enabled whenever \texttt{live} is non-empty, and
\texttt{Identify} is enabled for any \texttt{c? \notin live} whenever one
exists.  In the sole state where \texttt{Identify} has no fresh argument
--- \texttt{live = CONN} --- \texttt{Disconnect} is enabled.  So some
operation is always enabled.

% ============================================================
\section{Verification: the goals that must flip}
\label{sec:verify}
% ============================================================

Every command below names \emph{this} buggy file, not the correct spec ---
a reader who runs them against this variant sees the defect, which is the
whole point of a fidelity control.  The correct spec
(\texttt{crosshost\_hubid\_preemption.tex}) returns the opposite verdict on
each; its Fidelity section narrates the two traces.

Type-checking is a \texttt{fuzz} pass (this variant is well-typed --- the
defect is semantic, not syntactic):
\begin{verbatim}
  fuzz docs/crosshost_hubid_preemption_name_keyed_buggy.tex
\end{verbatim}

\paragraph{I1's negation must be \emph{found} (the double-owner).} Trace:
\texttt{Identify}(c1, h1, n1); \texttt{Identify}(c2, h1, n2) --- same
\texttt{HubId} \texttt{h1}, different name, so name-keyed preemption misses
\texttt{c1} and both stay live under \texttt{h1}.
\begin{verbatim}
  # must be FOUND against this variant (NOT found against the correct spec)
  probcli docs/crosshost_hubid_preemption_name_keyed_buggy.tex \
    -model_check -nodead \
    -goal "#(c1,c2).(c1 : live & c2 : live & not(c1 = c2) &
                     hubOf(c1) = hubOf(c2))" \
    -p DEFAULT_SETSIZE 3
\end{verbatim}

\paragraph{Witness W2 must be \emph{not found} (coexistence destroyed).}
Under name-keying a same-name arrival evicts the incumbent, so two live
connections can never share a name --- even with distinct \texttt{HubId}s,
both legitimate Hubs on different machines.
\begin{verbatim}
  # must be NOT found against this variant (FOUND against the correct spec)
  probcli docs/crosshost_hubid_preemption_name_keyed_buggy.tex \
    -model_check -nodead \
    -goal "#(c1,c2).(c1 : live & c2 : live & not(c1 = c2) &
                     not(hubOf(c1) = hubOf(c2)) &
                     nameOf(c1) = nameOf(c2))" \
    -p DEFAULT_SETSIZE 3
\end{verbatim}

Both discriminators come from the one key choice (\texttt{nameOf} in
\texttt{Identify}'s preempted set where the correct spec has
\texttt{hubOf}), which is why W11 is a single change and this one variant is
sufficient fidelity for Invariant 2. Deadlock-freedom is unaffected and
not the finding here; the correct spec model-checks it.

\end{document}
